\documentclass[11pt]{article}

% Revised working-paper manuscript.

\usepackage{authblk}
\usepackage{amsmath}
\usepackage{graphicx}
\graphicspath{{figures/}}

% Figures generated by the ecoChina2 visualization workflow are expected in
% the local figures/ folder. If a figure has not yet been copied from GitHub,
% the manuscript still compiles and shows the expected filename.
\newcommand{\projectfigure}[2][]{%
  \IfFileExists{figures/#2}{%
    \includegraphics[#1]{figures/#2}%
  }{%
    \fbox{\parbox{0.92\linewidth}{\centering
      Missing figure file: \texttt{\detokenize{#2}}\\
      Copy it from the GitHub folder
      \texttt{visualization var threshold0.4/figures/}.}}%
  }%
}
\usepackage{booktabs}
\usepackage[margin=1in]{geometry}
\usepackage{subcaption}
\captionsetup[table]{justification=centering,singlelinecheck=false}
\usepackage{threeparttable}
\usepackage{tabularx}
\usepackage{placeins}


% Citations / bibliography
\usepackage[authoryear,round]{natbib}
\usepackage[colorlinks=true,citecolor=blue,linkcolor=blue,urlcolor=blue]{hyperref}

\title{Finding Home for Forest Trees}

\author[1]{Jing Jiang}
\author[1]{Tongli Wang\thanks{Corresponding author: \texttt{tongli.wang@ubc.ca}}}

\affil[1]{Department of Forest and Conservation Sciences, Faculty of Forestry,
The University of British Columbia, Vancouver, BC, Canada}

\date{\today}

\begin{document}

\maketitle
\thispagestyle{plain}

% Abstract Section
\begin{abstract}
Vegetation-based ecological niche models may conflate climatic and edaphic
constraints and obscure within-species source-environment variation relevant
to seed transfer. We fitted one-vs-rest Random Forest models to the climatic
and topsoil associations of 53 ecotypes across China, gated projected climate
suitability by static soil suitability under SSP245 and SSP585 (2011--2100),
and used species--ecotype combinations as ecological proxies for source
populations of ten forest tree species. Climate models attained higher
held-out performance than topsoil models (mean TSS 92.9\% vs.\ 70.8--70.9\%).
Permutation-importance profiles were highly consistent between the two RF
workflows and revealed category-specific climatic and edaphic associations.
Exact reference-period Top-1 agreement was 60.2--60.3\%, but the observed
ecotype occurred within the five highest-ranked suitability analogues across
91.4\% of evaluated pixels. By 2071--2100 under SSP585, 23.9--24.2\% of mapped
area retained its reference-period Top-1 analogue, another current ecotype
became Top-1 across 64.8--65.1\%, and 10.7--11.3\% (0.91--0.96 million
km$^2$) had dual suitability below 0.4 for all 53 ecotypes. Among areas with a
different current Top-1 analogue, 76.3--76.4\% still retained the former
analogue within the future Top-5. Source-ecotype projections also diverged
within species: for \textit{Chamaecyparis formosensis}, a $\sim$7\% increase
in its species-niche envelope coincided with source-ecotype niche changes from
$-81\%$ to $+31\%$. Separating climatic from edaphic suitability, hard
assignments from ranked analogues, and species envelopes into source-ecotype
components identifies contrasts relevant to provenance matching and assisted
migration without predicting colonization, demographic persistence, or
ecosystem reassembly.
\end{abstract}

% Sections
\section{Introduction}

\citet{hutchinson1957} defined the $n$-dimensional hypervolume niche, where each axis in the multidimensional space represents an ecological factor. This foundational understanding of niche underscores the relationship between species and their ``home climate'' over evolutionary timescales. However, the escalating threat of climate change, marked by global surface temperature being 1.09~$^\circ$C higher in 2011--2020 than in 1850--1900 \citep{ipcc2021}, is challenging these established niches. Accompanied by more frequent extreme events such as heatwaves and severe drought, forest ecosystems are experiencing mismatches between their evolutionary ``home climate'' and rapidly changing conditions. This incompatibility is altering species ranges, with treelines moving poleward or to higher elevations \citep{hansson2021}. However, trees, with limited migratory capacities, face mortality risks when adaptation is slow and plasticity is constrained. Varied species responses to climate change disrupt established ecological interactions and introduce new dynamics \citep{lurgi2012,schweiger2008,schweiger2012}. Concurrently, altered climatic timelines result in phenological mismatches, in which seasonal events such as flowering become desynchronized \citep{bogdziewicz2021}.

For effective biodiversity conservation and sustainable forest management,
projecting environmental suitability for existing ecosystem classes is
important. Such projections can inform assisted migration, afforestation, and
reforestation planning. Climatic niche models correlate the spatial
distributions of classified ecosystems with reference-period climate and then
project where analogous climatic conditions may occur in the future
\citep{hamannWang2006,nwe2020,rehfeldt2012,wang2012,zhang2022}. Contraction in
projected suitable area indicates declining environmental analogy, whereas
expansion identifies potential recipient environments for assisted migration
\citep{oneill2008}; neither outcome directly predicts realized vegetation.
Despite similar methods, interpretations vary. \citet{nwe2020} describe their
models as predicting future ecosystem distributions, whereas
\citet{hamannWang2006}, \citet{wang2012}, and \citet{zhang2022} project the
future geographic distributions of climatic niches for existing ecosystem
classes. \citet{rehfeldt2012} discusses both realized climatic niches and
ecosystem distributions, illustrating this interpretive ambiguity.

Ecosystem-classification-based niche modeling has received several critiques. A particularly salient challenge is the concept that ecosystems can shift or migrate as a whole in response to changing climates. Paleoclimate evidence suggests that species' mobility responses to climate change differ, with each species moving or adapting at its own rate and in its own direction \citep{davisShaw2001}. Recent records also indicate that individual species within an ecosystem respond differently to climatic change; novel ecosystems may therefore be emerging \citep{lurgi2012,pecl2017}. In addition to differences in migration and adaptation among species, competition has been proposed as a challenge to suitability modeling for stable ecosystems. Previous models have also often overlooked soil. Soil, a crucial determinant of ecosystem type and health, encompasses complex interactions among minerals, organic matter, air, and water. It plays an integral role in biogeochemical cycles and supports plant growth through nutrient and water provision \citep{binkleyFisher2019}. Omitting soil may overestimate suitable habitat, whereas combining correlated soil and climate features in one model may underestimate future suitability. Some species-occurrence-based niche models have therefore modeled climate and soil separately and overlaid the two predictions to generate a suitability map \citep{feng2020}. More recently, \citet{wangCampbell2026} modeled ecosystem classes individually, integrated the class-specific RF predictions, and identified novel future climate space where no current ecosystem class exceeded a suitability threshold of 0.4.

The present framework treats projections as suitable environmental analogues
for current ecosystem classes, not forecasts of wholesale ecosystem migration.
Following \citet{wang2012}, it can inform assisted-migration planning without
assuming that all component species move together. It also identifies
locations lacking an adequate analogue among current classes, addressing the
novel-climate concern raised by \citet{rehfeldt2012}. Ecosystem classes remain
an imperfect but practical representation of biodiversity where finer data
are unavailable, and the separate soil component follows the overlay logic of
\citet{feng2020}.

Random Forest models \citep{breiman2001} combine multiple decision trees and provide a consensus estimate for ecosystem niche prediction. Through bagging and random predictor subsets \citep{amit1997,ho1998}, they can accommodate complex relationships among variables. Previous studies have used multiclass RFs to model ecosystem climate niches \citep{rehfeldt2012,robertsHamann2012,wang2012}. In contrast, this study uses individual-zone RFs, fitting each predefined ecotype separately to estimate its environmental suitability. Previous ecosystem niche models also often omitted soil or combined soil directly with climate variables \citep{nakazato2010}. Because soil and climate are correlated, combining them may obscure variable importance and the effect of climate change. Although some studies use variance inflation factors to reduce predictor collinearity \citep{namitha2022}, combining climate and soil in the same projection can still obscure their separate roles. Modeling climate and soil niches separately and then overlaying them provides a clearer interpretation. This approach has been used in species niche modeling but remains uncommon in ecosystem niche modeling \citep{feng2020}.

Individual one-vs-rest models also introduce a distinction between binary
performance and joint map assignment. A focal ecotype can be well separated
from sampled background environments while ranking below another ecotype at a
given cell. Retaining only the maximum suitability can therefore conceal
support for alternative environmental analogues. Examining ranked analogues
alongside the hard Top-1 assignment provides a direct diagnostic of niche
overlap and assignment sensitivity. Model interpretation further requires
identifying climatic and edaphic predictors that contribute consistently
within and among broad vegetation categories, while recognizing that RF
importance is predictive rather than causal.

The concept introduced by \citet{hutchinson1957} establishes the basis of this study, emphasizing that ecosystem classes occupy hypervolumes of environmental variables. Although established ecosystems represent relatively stable states resulting from extended succession, delineating potential new ecosystems would ideally require process-based models that simulate biotic and abiotic interactions over time. This study's correlative model instead provides a simplified representation linking stable states with climate and soil. Although it does not capture ecosystem dynamics, it may inform assisted migration of tree assemblages rather than only individual species. Competition and cooperation within and among tree species are important in ecosystems. For example, fir and birch exchange carbon and nitrogen through ectomycorrhizae \citep{simard1997}. Projecting suitable habitats for predefined ecosystems containing coexisting species can therefore inform multispecies afforestation decisions. Furthermore, where data on within-species population variation are scarce, ecotypes shaped by environmental variation may serve as ecological proxies for population differentiation \citep{wang2020,wangCampbell2026}. Projecting the niches of defined ecotypes can therefore inform in situ \citep{hamann2005} and ex situ \citep{wang2020} conservation and assisted migration using appropriate populations.

Owing to its extensive geographical and biological diversity, China offers a representative range of ecological variation for this study. Its climatic zones, ecosystems, and soil types provide a robust setting for testing the models. China's exposure to climate change also makes it a compelling region for examining potential shifts in ecosystem niches and developing resilience strategies.

We addressed four questions: (1) how accurately do individual climate and soil
models characterize the realized environmental associations of the 53
ecotypes, and which predictor groups contribute to those models; (2) how does
reference-period agreement change when the observed ecotype is evaluated as a
ranked analogue rather than only as the Top-1 assignment; (3) how do the rank,
top-ranked assigned area, and threshold-defined novelty of current ecotype
niches change with time and forcing; and (4) how strongly do source-ecotype
population proxies diverge within species? These projections identify
environmental analogues relevant to provenance matching and assisted
migration, but do not evaluate colonization, demographic persistence, biotic
interactions, or the broader risks of moving organisms among ecosystems.

\section{Methods}

\subsection{Data sources and preparation}

\paragraph{Vegetation ecotypes.}
The vegetation layer was derived from the digital \textit{Vegetation Map of China} (1:1,000,000) provided by the Environmental \& Ecological Science Data Center for West China, National Natural Science Foundation of China. The third-level classification was used. The map was rasterized to 30-arc-second resolution in WGS84. Of the 56 mapped zones, 53 ecotypes with sufficient observations and valid definitions were retained for modeling. Each retained ecotype was subsequently modeled as a focal zone against all other retained zones.

Ecotype identifiers, names, and broad vegetation categories were retained in
the analysis lookup table distributed with the model outputs.

\paragraph{Climate.}
Reference-period climate (1961--1990) was obtained from ClimateAP through the ClimateNAr R package v2.0.0 \citep{wang2017}. Pixel-center coordinates and elevation from an 800-m digital elevation model were used to retrieve annual and seasonal climate variables. Future climate was obtained from the eight-GCM ensemble for SSP245 and SSP585 for 2011--2040, 2041--2070, and 2071--2100.

\paragraph{Soil.}
Topsoil data were obtained from the Harmonized World Soil Database (HWSD) \citep{fischer2012hwsd} at 30-arc-second resolution. Because subsoil coverage was incomplete in western China, the analysis used 15 topsoil variables (0--30 cm; Table~\ref{tab:topsoil-vars}). Soil predictors were treated as static in future projections.

\begin{table}[htbp]
  \centering
  \caption{Topsoil variables used in model building.}
  \label{tab:topsoil-vars}
  \small
  \setlength{\tabcolsep}{5pt}
  \renewcommand{\arraystretch}{1.08}
  \begin{tabularx}{\textwidth}{@{}>{\raggedright\arraybackslash}p{0.27\textwidth} >{\raggedright\arraybackslash}X@{}}
    \toprule
    Code & Variable / definition \\
    \midrule
    T\_GRAVEL & Gravel content (\% by volume) \\
    T\_SAND & Sand fraction (\%) \\
    T\_SILT & Silt fraction (\%) \\
    T\_CLAY & Clay fraction (\%) \\
    T\_REF\_BULK\_DENSITY & Bulk density (soil mass per unit volume) \\
    T\_OC & Organic carbon (\%) \\
    T\_PH\_H2O & pH measured in a soil--water solution \\
    T\_CEC\_CLAY & Cation exchange capacity of the clay fraction \\
    T\_CEC\_SOIL & Cation exchange capacity of topsoil \\
    T\_BS & Base saturation \\
    T\_TEB & Total exchangeable bases \\
    T\_CACO3 & Calcium carbonate content \\
    T\_CASO4 & Gypsum (calcium sulphate) content \\
    T\_ESP & Exchangeable sodium percentage \\
    T\_ECE & Electrical conductivity (salinity) \\
    \bottomrule
  \end{tabularx}
\end{table}

\paragraph{Tree species.}
Species occurrence layers for ten major forest trees were overlaid with the reference ecotype raster following \citet{wang2020}. Consistent with the ecosystem-class interpretation in \citet{wangCampbell2026}, ecotype membership was treated as an ecological proxy for within-species population differentiation. Within each species, a species--ecotype combination was treated as a population and retained when the species occupied at least 10 raster cells in that ecotype. Reference occupancy was the number of occupied cells. Shannon $H$ summarized the relative occupancy distribution among retained populations within each species. This produced 89 reference populations across the ten species.

\begin{table}[htbp]
  \centering
  \caption{Tree species, population counts, and population diversity.}
  \label{tab:tree-pops}
  \scriptsize
  \setlength{\tabcolsep}{3.2pt}
  \renewcommand{\arraystretch}{1.08}
  \begin{tabularx}{\textwidth}{@{}l>{\raggedright\arraybackslash}Xrrrr@{}}
    \toprule
    Code & Species &
    \shortstack{Reference\\populations} &
    \shortstack{Reference\\Shannon $H$} &
    \shortstack{Projection-eligible\\populations} &
    \shortstack{Eligible\\Shannon $H$} \\
    \midrule
    querVar & \textit{Quercus variabilis} & 18 & 1.505 & 17 & 1.389 \\
    pinuYun & \textit{Pinus yunnanensis} & 16 & 0.385 & 15 & 0.360 \\
    phylPub & \textit{Phyllostachys pubescens} & 9 & 0.455 & 8 & 0.406 \\
    lariGme & \textit{Larix gmelinii} & 9 & 0.201 & 8 & 0.180 \\
    robiPse & \textit{Robinia pseudoacacia} & 8 & 0.501 & 7 & 0.452 \\
    lariOlg & \textit{Larix olgensis} var.~\textit{changpaiensis} & 8 & 0.407 & 7 & 0.362 \\
    chamFor & \textit{Chamaecyparis formosensis} & 7 & 1.818 & 7 & 1.818 \\
    cyclLon & \textit{Cyclobalanopsis longinux} & 6 & 1.155 & 6 & 1.155 \\
    saliMat & \textit{Salix matsudana} & 4 & 0.615 & 3 & 0.410 \\
    pinuSyl & \textit{Pinus sylvestris} var.~\textit{mongolica} & 4 & 0.164 & 4 & 0.164 \\
    \bottomrule
  \end{tabularx}
  \begin{minipage}{0.98\textwidth}
    \footnotesize
    \vspace{0.3em}
    \textbf{Notes.} Shannon $H$ summarizes the relative occupancy distribution among retained
    populations within each species. Projection-eligible populations had the
    required source-ecotype dual-suitability outputs.
  \end{minipage}
\end{table}

\subsection{Individual-zone random forest modeling and evaluation}

Each ecotype was modeled independently in a one-vs-rest classification. Let $\mathcal{Z}$ be the set of 53 modeled ecotypes and let $G(i)\in\mathcal{Z}$ be the reference ecotype at raster cell $i$. For focal ecotype $z\in\mathcal{Z}$, the response was
\begin{equation}
Y_z(i)=
\begin{cases}
1, & G(i)=z,\\
0, & G(i)\in\mathcal{Z}\setminus\{z\},
\end{cases}
\end{equation}
where $Y_z(i)$ is the binary response. Separate models produced climate suitability $S_{\mathrm{clim},z}(i)$ and soil suitability $S_{\mathrm{soil},z}(i)$.

After removing missing observations, climate and soil data were split 70:30 into training and held-out test sets with a fixed random seed. Focal-ecotype cells were presences; cells from all other modeled ecotypes were absences. Ecotype-specific predictors were retained by backward-purging selection. We fitted two final workflows: a single 500-tree RF (\texttt{rf\_var}) and a multi-Forest RF (\texttt{mf\_var}) that aggregated ten separately grown 100-tree forests. Training samples contained at most 8,000 presences, 1.3 absences per presence, and 20,000 observations in total. Because the sampled absence pool also contained approximately 1.3 absences per presence, each component forest used the same sampled training pool. The multi-Forest workflow therefore tested sensitivity to tree number and forest aggregation, not to repeated absence sampling.

Held-out evaluation used balanced subsets containing all available focal-ecotype presences and an equal number of absences. We report balanced accuracy, sensitivity, specificity, precision, F1, the true skill statistic (TSS), and area under the receiver operating characteristic curve (AUC). AUC was calculated from class-1 vote probabilities using the rank-based Mann--Whitney estimator; other metrics used a 0.5 threshold. Reference-period assigned maps were compared cellwise with the reference ecotype raster using coverage, exact-zone agreement, broad-category agreement, and ecotype-level one-vs-rest confusion metrics. AUC was not calculated for this hard-label overlay.

For comparison, we fitted a multiclass RF with ecotype as the categorical response and both climate and soil predictors. Sampling within each tree was balanced by class. We compared its reference-period map with the individual-zone dual-niche maps using the same map-level metrics. This model was used only to assess robustness to model formulation, not for future, population, or species projections.

\subsection{Feature importance and ecological summaries}

Feature importance was extracted from the saved climate and soil RF objects
without refitting the models. The primary measure was unscaled global
permutation importance (\texttt{MeanDecrease\allowbreak Accuracy}) from
\texttt{randomForest}. Raw values were retained for auditing; positive values
were normalized to sum to one within each ecotype, niche component, and
workflow, whereas negative values contributed zero to normalized shares.
\texttt{MeanDecreaseGini} was retained only as a secondary diagnostic.

Plain RF and Plain MF RF importance shares were averaged within each ecotype
to form a consensus profile. Predictors were then summarized by climatic or
edaphic factor group and, for climate, by annual or seasonal period. Category
summaries are unweighted means with SE across ecotypes, so large ecotypes do
not dominate the inferred category profile. Agreement between RF workflows
was evaluated using Pearson and Spearman correlations of the complete
importance profiles and Jaccard overlap among the five highest-ranked
predictors. Descriptive within-category comparisons used the profile medoid,
defined as the ecotype with the highest mean pairwise cosine similarity, and
the ecotypes with the largest and smallest observed reference areas.
Reader-facing summaries excluded the no-vegetation class, although it remained
in completeness audits. Permutation importance identifies contributions to
prediction, not response direction, response shape, mechanism, or causality;
interpretation therefore emphasized recurrent factor groups and agreement
across workflows rather than individual ranks among correlated predictors.

\subsection{Ecotype niche projection, assignment, and changes in suitable area}

Class-1 RF probabilities were projected with \texttt{terra::predict}. Climate suitability was projected for 1961--1990 and three future periods (2011--2040, 2041--2070, and 2071--2100) under SSP245 and SSP585; soil suitability was projected once and held constant. Let $S_{\mathrm{clim},z,t,s}(i)$ and $S_{\mathrm{soil},z}(i)$ denote climate and soil suitability, respectively, for ecotype $z$, raster cell $i$, period $t$, and scenario $s\in\{\mathrm{SSP245},\mathrm{SSP585}\}$; the scenario subscript is omitted for the reference period.

Climate and soil suitability were combined using a fixed soil gate:
\begin{equation}
S_{\mathrm{dual},z,t,s}(i)=
\begin{cases}
S_{\mathrm{clim},z,t,s}(i), & S_{\mathrm{soil},z}(i)>0.2,\\
0, & S_{\mathrm{soil},z}(i)\leq 0.2,\\
\mathrm{NA}, & S_{\mathrm{soil},z}(i)=\mathrm{NA}.
\end{cases}
\end{equation}
Thus, soil acted as a static gate at a strict threshold of 0.2, and missing soil suitability remained NA.

Each cell was assigned to the ecotype with the highest dual suitability:
\begin{equation}
z^*_{t,s}(i)=\underset{z}{\operatorname{arg\,max}}\;S_{\mathrm{dual},z,t,s}(i),
\end{equation}
where $z^*_{t,s}(i)$ is the top-ranked ecotype-niche assignment. Suitability
values within $10^{-4}$ of the maximum were treated as ties. If the
reference-map ecotype was among the tied candidates, it was retained;
otherwise, one candidate was selected using a fixed random seed. Following
\citet{wangCampbell2026}, the 0.4 cutoff was applied to dual suitability:
future cells with $\max_z S_{\mathrm{dual},z,t,s}(i)<0.4$ were assigned to
novel niche space (Zone 99), which was excluded from reference-period
assignments. This definition depends only on the absolute dual-suitability
threshold and is independent of Top-$k$ rank.

Positive dual-suitability values were ranked within each cell, and the ten
highest-ranked ecotypes and values were retained. Reference-period Top-$k$
agreement was evaluated for $k\in\{1,2,3,5\}$. Top-1 agreement compared the
assigned map with the reference ecotype map and preserved the tie rule. For
$k>1$, agreement indicated that the reference ecotype occurred among the first
$k$ ranked analogues. Results were calculated as both pixel shares and
cell-area-weighted shares; pixel shares were used when comparing Top-$k$
agreement with exact map agreement.

The primary future rank analysis used each workflow's own reference-period
Top-1 assigned ecotype as a fixed baseline. On a common valid mask, cells were
classified into five mutually exclusive outcomes: the reference-period
analogue remained Top-1; another current ecotype became Top-1 while the former
analogue ranked second or third; the former analogue ranked fourth or fifth;
the former analogue fell below the Top-5; or no current ecotype reached the
0.4 threshold. Conditional Top-3 and Top-5 retention was also calculated among
cells assigned to another current ecotype. This analysis evaluates changes in
relative niche rank, except for the threshold-defined novel class; membership
in the Top-3 or Top-5 is not an additional claim of adequate suitability.

Two supplementary sensitivity analyses were retained. Analogue richness was
defined as the number of current ecotypes with dual suitability of at least
0.4; areas with zero, fewer than three, and fewer than five such analogues were
summarized separately. A matched Top-$k$ analysis applied the same
reference-conditioned recovery rule to reference and future maps: for
$k\in\{3,5\}$, the observed ecotype label was restored when it occurred within
the first $k$ ranks, otherwise the period-specific Top-1 assignment was
retained. Because this construction uses the observed label, it was treated as
a sensitivity analysis rather than an independent future projection.

Changes in top-ranked assignment were measured from each workflow's
reference-period assigned map, not from the reference vegetation raster. All
assignment, novel-space, population, and species areas were calculated using
cell-specific values from \texttt{terra::cellSize} rather than raster-cell
counts. These quantities describe modeled niche suitability and rank, not
realized ecosystem composition.


\subsection{Tree species niche projection}

Population $p$ inherited the dual suitability of its source ecotype $z_p$:
\begin{equation}
S_{\mathrm{pop},p,t,s}(i)=S_{\mathrm{dual},z_p,t,s}(i),
\end{equation}
where $S_{\mathrm{pop},p,t,s}(i)$ is population suitability and $i$, $t$, and $s$ denote cell, period, and scenario. Suitable area summed cell-specific area where dual suitability was at least 0.4; suitability-weighted area summed cell-specific area multiplied by dual suitability.

For species $q$, population niches in the retained set $P_q$ were combined by their cellwise maximum:
\begin{equation}
S_{\mathrm{species},q,t,s}(i)=\max_{p\in P_q}S_{\mathrm{pop},p,t,s}(i),
\end{equation}
where $S_{\mathrm{species},q,t,s}(i)$ is species-level suitability. When multiple populations were suitable in a cell, their source ecotypes and corresponding dual-suitability values were retained in ranked order.



\section{Results}

\subsection{Predictive performance and environmental structure of ecotype niches}

Across the 53 ecotypes, climate models attained higher held-out performance
than topsoil models (Fig.~\ref{fig:rf_performance};
Table~\ref{tab:rf_overall_final}). The two RF workflows gave nearly identical
estimates. Climate-model balanced accuracy was 96.43--96.44\%, F1 was
96.50--96.51\%, TSS was 92.86--92.88\%, and mean AUC was
0.98848--0.98855. Corresponding soil-model values were 85.42--85.43\%,
86.03--86.05\%, 70.84--70.87\%, and 0.87986--0.88186. Soil-model sensitivity
(89.63--89.66\%) exceeded specificity (81.21\%), and between-ecotype
variation was greater for soil than climate models.

\begin{table}[!htbp]
\centering
\caption{Mean (SE) held-out performance across 53 ecotypes.}
\label{tab:rf_overall_final}
\begin{threeparttable}
\centering
\scriptsize
\setlength{\tabcolsep}{2.0pt}
\begin{tabular}{@{}llrrrrrrr@{}}
\toprule
Niche & Workflow &
\shortstack{Balanced\\accuracy (\%)} &
\shortstack{Sensitivity\\(\%)} &
\shortstack{Specificity\\(\%)} &
\shortstack{Precision\\(\%)} &
F1 (\%) &
TSS (\%) &
AUC \\
\midrule
Climate & Plain RF    & 96.43 (0.32) & 97.90 (0.23) & 94.96 (0.44) & 95.15 (0.41) & 96.50 (0.31) & 92.86 (0.64) & 0.9885 (0.0014) \\
Climate & Plain MF RF & 96.44 (0.32) & 97.91 (0.23) & 94.97 (0.44) & 95.16 (0.41) & 96.51 (0.31) & 92.88 (0.64) & 0.9885 (0.0014) \\
\addlinespace
Soil & Plain RF       & 85.42 (0.65) & 89.63 (0.70) & 81.21 (0.93) & 82.88 (0.74) & 86.03 (0.62) & 70.84 (1.31) & 0.8799 (0.0058) \\
Soil & Plain MF RF    & 85.43 (0.65) & 89.66 (0.70) & 81.21 (0.92) & 82.89 (0.73) & 86.05 (0.62) & 70.87 (1.30) & 0.8819 (0.0058) \\
\bottomrule
\end{tabular}
\begin{tablenotes}[flushleft]
\footnotesize
\item[] \textbf{Notes.} Entries are means with SE in parentheses across 53
ecotypes; held-out test sets used a 1:1 presence--absence ratio. All metrics
except AUC are percentages. AUC is the rank-based Mann--Whitney estimator and
was finite for all ecotypes. Sensitivity is the true-positive rate;
specificity is the true-negative rate; precision is the positive predictive
value; TSS is sensitivity + specificity - 1. MF denotes multi-Forest.
\end{tablenotes}
\end{threeparttable}
\end{table}

\begin{figure}[!htbp]
  \centering
  % GitHub: visualization var threshold0.4/figures/Figure_var_1_climate_soil_model_performance.png
  \projectfigure[width=0.98\linewidth]{Figure_var_1_climate_soil_model_performance.png}
  \caption{Held-out performance of the final climate and soil RF models.
  Points show means across 53 ecotypes and error bars show SE.}
  \label{fig:rf_performance}
\end{figure}

Permutation-importance profiles were highly consistent between Plain RF and
Plain MF RF. After excluding the no-vegetation class from ecological summaries,
mean Spearman correlations across 52 ecotypes were 0.987 for climate and 0.975
for soil; median Top-5 Jaccard overlap was 1.0 for both niche components. The
shared training samples and predictor sets therefore produced nearly identical
importance rankings under the two forest configurations.

Climatic importance nevertheless differed among vegetation categories
(Fig.~\ref{fig:feature_importance_groups}a). Temperature accounted for the
largest mean share in alpine vegetation (31.4\%), forest (27.5\%),
grassland--meadow--steppe (24.9\%), and scrub (36.6\%). Growing-season and
degree-day variables ranked first in cropland (32.5\%) and desert (24.6\%),
whereas precipitation ranked first in wetland (23.0\%). Annual variables formed
the largest single temporal group in each category (28.8--36.2\%), but seasonal
variables together accounted for 63.8--71.2\% of climatic importance.

Nutrient retention and base status dominated topsoil importance in cropland,
forest, grassland--meadow--steppe, scrub, and wetland (40.4--46.8\%;
Fig.~\ref{fig:feature_importance_groups}b). Alpine vegetation had a more even
edaphic profile, with nutrient retention and base status (25.8\%), coarse
fragments (21.8\%), and texture (18.9\%) all prominent. Desert ecotypes also
differed from the general pattern, with contributions from organic carbon
(17.9\%) alongside nutrient retention and base status (22.6\%).

Within-category contrasts showed that these means did not erase ecotype-level
variation. For example, the large alpine sparse-vegetation ecotype (Zone 20)
emphasized summer evaporative demand, summer maximum temperature, and coarse
fragments, whereas the much smaller alpine-tundra ecotype (Zone 17) emphasized
winter maximum temperature, annual evaporative demand, organic carbon, and
base status. Within forest, the large subtropical needleleaf ecotype (Zone 30)
emphasized growing-season heat and precipitation, whereas the smaller
subtropical montane mixed-forest ecotype (Zone 34) emphasized temperature
range, spring precipitation, and edaphic texture or base status.

\begin{figure}[!htbp]
  \centering
  \begin{subfigure}{0.98\linewidth}
    \centering
    % GitHub: assessment_var/feature_importance_analysis/reader_figures/Figure_FI_reader_category_factor_groups_climate.png
    \projectfigure[width=\linewidth]{Figure_FI_reader_category_factor_groups_climate.png}
    \caption{Climatic factor groups.}
  \end{subfigure}

  \vspace{0.5em}

  \begin{subfigure}{0.98\linewidth}
    \centering
    % GitHub: assessment_var/feature_importance_analysis/reader_figures/Figure_FI_reader_category_factor_groups_soil.png
    \projectfigure[width=\linewidth]{Figure_FI_reader_category_factor_groups_soil.png}
    \caption{Topsoil factor groups.}
  \end{subfigure}
  \caption{Category-level shares of normalized permutation importance.
  Plain RF and Plain MF RF profiles were averaged within ecotype before
  calculating unweighted category means. Importance describes predictive
  contribution, not effect direction or causality.}
  \label{fig:feature_importance_groups}
\end{figure}

\FloatBarrier

\subsection{Reference-period assignment and ranked niche analogues}

When the 53 dual-niche models were compared to assign one top-ranked analogue
per cell, exact agreement with the reference ecotype map was lower than the
individual one-vs-rest metrics. Both workflows covered 98.15\% of valid
reference cells. Exact-zone agreement was 60.27\% for Plain RF and 60.24\% for
Plain MF RF; broad-category agreement was 69.94\% for both. On the common valid
mask, macro ecotype-level balanced accuracy was 82.49--82.51\%, F1 was
49.84--49.87\%, and TSS was 64.99--65.01\%
(Table~\ref{tab:reference_map_final}).

\begin{table}[!htbp]
\centering
\caption{Reference-period map agreement of the final individual-zone RF workflows.}
\label{tab:reference_map_final}
\begin{threeparttable}
\centering
\scriptsize

\textit{Map agreement}\\[-0.2em]
\setlength{\tabcolsep}{6pt}
\begin{tabular}{@{}lrrr@{}}
\toprule
Workflow & Coverage (\%) & Exact zone (\%) & Broad category (\%) \\
\midrule
Plain RF    & 98.15 & 60.27 & 69.94 \\
Plain MF RF & 98.15 & 60.24 & 69.94 \\
\bottomrule
\end{tabular}

\vspace{0.6em}

\textit{Macro ecotype-level performance}\\[-0.2em]
\setlength{\tabcolsep}{2.7pt}
\begin{tabular}{@{}lrrrrrr@{}}
\toprule
Workflow &
\shortstack{Balanced\\accuracy (\%)} &
\shortstack{Sensitivity\\(\%)} &
\shortstack{Specificity\\(\%)} &
\shortstack{Precision\\(\%)} &
F1 (\%) &
TSS (\%) \\
\midrule
Plain RF    & 82.49 (1.11) & 65.76 (2.17) & 99.23 (0.06) & 45.40 (2.86) & 49.87 (2.18) & 64.99 (2.22) \\
Plain MF RF & 82.51 (1.11) & 65.78 (2.17) & 99.23 (0.06) & 45.37 (2.86) & 49.84 (2.18) & 65.01 (2.21) \\
\bottomrule
\end{tabular}

\begin{tablenotes}[flushleft]
\footnotesize
\item[] \textbf{Notes.} Coverage and agreement are map-wide percentages.
Macro entries are means with SE in parentheses across 53 ecotype-level
one-vs-rest confusion matrices on the common valid mask. AUC is not defined
for the hard-label assigned-map overlay.
\end{tablenotes}
\end{threeparttable}
\end{table}

Ranked agreement showed that most Top-1 disagreements did not place the
reference ecotype far outside the leading suitability candidates
(Fig.~\ref{fig:reference_topk}). Across 12,393,826 evaluated pixels, inclusion
of the reference ecotype increased from 60.24--60.27\% at Top-1 to
86.48--86.49\% at Top-3 and 91.35\% at Top-5. The corresponding area-weighted
shares were 59.86--59.88\%, 86.22--86.23\%, and 91.28\%. Among pixels without
exact Top-1 agreement, 78.2\% still placed the reference ecotype within the
Top-5. The remaining 8.6\% of evaluated pixels placed it outside the Top-5.

\begin{figure}[!htbp]
  \centering
  \begin{subfigure}{0.98\linewidth}
    \centering
    % GitHub: visualization var threshold0.4/figures/Reference_maps_plain_rf_plain_mf.png
    \projectfigure[width=\linewidth]{Reference_maps_plain_rf_plain_mf.png}
    \caption{Predicted reference-period assigned maps.}
  \end{subfigure}

  \vspace{0.5em}

  \begin{subfigure}{0.94\linewidth}
    \centering
    % GitHub: visualization var threshold0.4/figures/Figure_var_4_reference_topk_agreement_decomposition.png
    \projectfigure[width=\linewidth]{Figure_var_4_reference_topk_agreement_decomposition.png}
    \caption{First rank interval in which the reference ecotype was recovered.}
  \end{subfigure}
  \caption{Reference-period assignment and ranked analogue agreement. Top-1
  is exact agreement under the assigned-map tie rule; Top-3 and Top-5 indicate
  whether the reference ecotype occurred among the first $k$ ranked
  dual-suitability analogues. Panel (b) reports pixel shares.}
  \label{fig:reference_topk}
\end{figure}

\FloatBarrier

\subsection{Future retention of current niche analogues and novel niche space}

Projected changes were measured relative to each workflow's reference-period
Top-1 assigned map. Under Plain RF and SSP245, the reference-period analogue
remained Top-1 across 52.58\% of mapped area in 2011--2040, another current
ecotype became Top-1 across 45.47\%, and 1.95\% was novel. By 2071--2100, the
retained Top-1 share declined to 35.01\% under SSP245 and 24.24\% under SSP585,
whereas the different-current-analogue shares increased to 59.27\% and 65.07\%,
respectively (Figs.~\ref{fig:future_niche_maps} and
\ref{fig:future_rank_novel}). Plain MF RF gave corresponding late-century
SSP585 shares of 23.88\%, 64.78\%, and 11.34\%.

The former analogue often remained competitive after losing the Top-1 rank.
Among areas assigned to another current ecotype in 2011--2040 SSP245,
73.02--73.26\% retained the former analogue within the future Top-3 and
92.54--92.70\% retained it within the Top-5. By 2071--2100 SSP585, these shares
declined to 46.69--47.26\% and 76.32--76.35\%. Across all evaluated area in the
latter period, 23.88--24.24\% remained Top-1, 30.38--30.61\% placed the former
analogue at ranks 2--3, 18.80--19.28\% placed it at ranks 4--5,
15.34--15.39\% placed it below the Top-5, and 10.70--11.34\% was novel.

Novel niche-space area increased with time and forcing. For Plain RF under SSP245, it was
165,900~km$^2$ (1.95\%) in 2011--2040, 324,100~km$^2$ (3.81\%) in
2041--2070, and 486,400~km$^2$ (5.72\%) in 2071--2100. Under SSP585, it was
174,000~km$^2$ (2.05\%) initially, 471,900~km$^2$ (5.55\%) by mid-century,
and 909,400~km$^2$ (10.70\%) by late century. Plain MF RF gave similar
late-century estimates: 477,500~km$^2$ (5.62\%) under SSP245 and
964,400~km$^2$ (11.34\%) under SSP585.

Analogue richness provided a graded complement to Zone 99. Under Plain RF and
SSP585, the area with fewer than three current ecotypes reaching dual
suitability 0.4 increased from 37.2\% in 2011--2040 to 61.4\% in 2071--2100;
the area with fewer than five increased from 77.4\% to 91.1\%. These nested
criteria measure the size of the suitable-analogue set, not novel area.

Top-ranked assigned area also differed among ecotypes. From the reference
period to 2071--2100 SSP585 under Plain RF, Zone 21, a managed
grain-field--orchard--plantation mosaic, increased by 944,000~km$^2$; Zones 27
and 18 increased by 346,900 and 320,500~km$^2$. The largest decreases occurred
for Zone 16 high-cold meadow ($-333{,}900$~km$^2$), Zone 30 subtropical
needleleaf forest ($-298{,}500$~km$^2$), and Zone 31 high-cold steppe
($-248{,}000$~km$^2$; Supplementary Table~\ref{tab:zone_assignment_change_si}).
These are changes in the area where a niche was ranked first, not projections
of realized ecosystem extent.

\begin{figure}[!htbp]
  \centering
  % GitHub: visualization var threshold0.4/figures/Figure_var_5_future_maps_rf_var.png
  \projectfigure[width=0.98\linewidth]{Figure_var_5_future_maps_rf_var.png}
  \caption{Projected Top-1 ecotype-niche assignments from Plain RF under
  SSP245 and SSP585. Zone 99 denotes novel niche space in which no current
  ecotype reaches dual suitability 0.4. The maps show modeled environmental
  analogues, not future vegetation composition.}
  \label{fig:future_niche_maps}
\end{figure}

\begin{figure}[!htbp]
  \centering
  \begin{subfigure}{0.98\linewidth}
    \centering
    % GitHub: visualization var threshold0.4/figures/Figure_var_11b_future_topk_retention.png
    \projectfigure[width=\linewidth]{Figure_var_11b_future_topk_retention.png}
    \caption{Future rank of each workflow's reference-period Top-1 analogue.}
  \end{subfigure}

  \vspace{0.5em}

  \begin{subfigure}{0.90\linewidth}
    \centering
    % GitHub: visualization var threshold0.4/figures/Figure_var_6_novel_ecosystem_area.png
    \projectfigure[width=\linewidth]{Figure_var_6_novel_ecosystem_area.png}
    \caption{Threshold-defined novel niche-space area.}
  \end{subfigure}
  \caption{Future displacement of current niche analogues and emergence of
  novel niche space. Panel (a) uses each workflow's reference-period Top-1
  assigned map as the baseline and partitions mapped area by the future rank of
  that analogue. Panel (b) reports cells for which all 53 ecotypes have dual
  suitability below 0.4; this definition is independent of Top-$k$. Areas are
  weighted by cell-specific surface area.}
  \label{fig:future_rank_novel}
\end{figure}

\FloatBarrier

\subsection{Source-ecotype population and species-niche contrasts}

The ten species comprised 89 reference species--ecotype populations, of which
82 had the dual-suitability outputs required for projection
(Table~\ref{tab:tree-pops}). \textit{Quercus variabilis} occurred in the most
reference ecotypes (18), followed by \textit{Pinus yunnanensis} (16).

Population count did not describe how occurrences were distributed among
source ecotypes. \textit{P. yunnanensis} had 16 populations but a reference
Shannon $H$ of 0.385, indicating concentrated occupancy, whereas
\textit{Chamaecyparis formosensis} had seven populations and the highest
Shannon $H$ (1.818), indicating a more even source-ecotype representation.

Under Plain RF and SSP585, changes in species-level area with dual suitability
of at least 0.4 from 2011--2040 to 2071--2100 ranged from $-31.3\%$ for
\textit{Larix gmelinii} to $+12.0\%$ for \textit{Cyclobalanopsis longinux}.
The retained occurrences of \textit{L. gmelinii} were concentrated in Zone 2,
the cold-temperate and montane-temperate needleleaf-forest ecotype; the largest
source population of \textit{C. longinux} was in Zone 49, the subtropical
broadleaf evergreen-forest ecotype. Several other species-level niches
decreased by 15--19\%, whereas that of \textit{C. formosensis} increased by
approximately 7\% (Fig.~\ref{fig:population_species_future}). These values
describe cellwise maxima across the source-ecotype niches retained for each
species, not projected species distributions.

\textit{C. formosensis} showed how species-level aggregation concealed
source-environment contrasts
(Fig.~\ref{fig:chamfor_population_detail}). Its seven source ecotypes included
montane or subtropical needleleaf and mixed forests (Zones 10, 30, and 34),
bamboo forest and scrub (Zone 26), subtropical broadleaf evergreen forest
(Zone 49), and managed subtropical or tropical crop--orchard--plantation
mosaics (Zones 21 and 53). Under Plain RF and SSP585, its species-level suitable
area increased from 2.55 to 2.73 million~km$^2$ ($\sim7\%$) between 2011--2040
and 2071--2100. Over the same interval, suitable area decreased by 81.4\% for
the Zone-26 population proxy, 53.5\% for Zone 30, and 40.8\% for Zone 49, but
increased by 20.3\% for Zone 21 and 30.5\% for Zone 53. The range between the
largest decrease and increase exceeded 110 percentage points.

Plain MF RF reproduced the principal contrast: Zone-26 and Zone-49 population
proxies decreased by 82.7\% and 41.7\%, whereas Zones 21 and 53 increased by
20.3\% and 32.6\%.

These associations do not imply that \textit{C. formosensis} is adapted to
cropland or plantation land cover; each label describes the composite ecotype
from which the population proxy was defined. For any given source ecotype, the
same dual-suitability surface was inherited by every species occurring there.
Species identity determined which source-ecotype surfaces were included in the
species-level maximum, not the values within those surfaces.

\begin{figure}[!htbp]
  \centering
  % GitHub: visualization var threshold0.4/figures/Figure_var_10a_population_species_summary.png
  \begin{subfigure}{0.96\linewidth}
    \centering
    \projectfigure[width=\linewidth]{Figure_var_10a_population_species_summary.png}
    \caption{Source-ecotype composition across focal species.}
  \end{subfigure}

  \vspace{0.6em}

  % GitHub: visualization var threshold0.4/figures/Figure_var_9_dual_species_area.png
  \begin{subfigure}{0.96\linewidth}
    \centering
    \projectfigure[width=\linewidth]{Figure_var_9_dual_species_area.png}
    \caption{Species-niche area with dual suitability $\geq 0.4$.}
  \end{subfigure}
  \caption{Source-ecotype population structure and projected species-level niche suitability. Species-level suitability is the cellwise maximum across retained source-ecotype population niches; area is calculated from cell-specific surface area.}
  \label{fig:population_species_future}
\end{figure}

\begin{figure}[!htbp]
  \centering
  % Copy from GitHub population-detail folder into the local figures/ folder.
  \projectfigure[width=0.96\linewidth]{Figure_var_10a_population_area_chamFor.png}
  \caption{Projected suitable-area trajectories for \textit{Chamaecyparis formosensis} source-ecotype population niches. Each line is the dual-suitability footprint inherited from one source ecotype; it is an environmental proxy, not an independently fitted population model.}
  \label{fig:chamfor_population_detail}
\end{figure}

\FloatBarrier

\subsection{Robustness across RF formulations}

On the common valid mask, exact-zone agreement was 60.34\% for Plain RF,
60.32\% for Plain MF RF, and 57.99\% for multiclass RF; broad-category
agreement was 70.02\%, 70.01\%, and 68.86\%, respectively. The multiclass RF
had higher macro sensitivity and TSS but lower exact agreement and F1 than the
individual-zone workflows (Supplementary Fig.~\ref{fig:robustness_multiclass_si}).
Plain RF and Plain MF RF were nearly indistinguishable in held-out metrics,
reference Top-$k$ agreement, future rank retention, novel area, population
projections, and feature-importance profiles. Their agreement supports
numerical stability to tree number and forest aggregation, but not to training
sample selection because they used the same sampled pool.

\FloatBarrier

\section{Discussion}

The models characterize environmental associations of reference-period ecotype
occurrences and therefore approximate realized niches within the selected
climate--soil predictor space, not full Hutchinsonian fundamental niches. Four
findings are central. Climate models attained higher held-out performance than
topsoil models, while feature importance revealed ecologically structured
variation within both components. Ranked analogues reconciled strong binary
performance with lower hard-map agreement. Future projections showed a
progressive decline in the rank of many reference-period analogues, distinct
from the smaller area of threshold-defined novelty. Finally, source-ecotype
population proxies exposed within-species contrasts hidden by species-level
envelopes.

\subsection{Climatic and edaphic structure of modeled ecological niches}

Higher climate-model performance indicates stronger separation in the
climatic predictor space represented here; it does not establish that climate
is universally more important than soil. Shared topsoil properties and omitted
controls---including hydrology, topography, substrate heterogeneity,
disturbance, and land use---may contribute to the lower and more variable soil
performance. The use of topsoil alone may be especially limiting where rooting
depth, groundwater, permafrost, or subsoil texture constrain vegetation.

The feature-importance profiles provide ecological context for this contrast.
Temperature, precipitation, and growing-season variables contributed most
across many categories, but their relative contributions followed recognizable
environmental differences. Alpine niches emphasized thermal and evaporative
conditions, cropland and desert niches emphasized growing-season variables,
and wetland niches placed greater weight on precipitation. Seasonal variables
collectively contributed more than annual variables, supporting the use of
seasonal climate descriptors for separating realized ecotype niches.

Topsoil importance was more strongly organized around nutrient retention and
base status, particularly in forest, cropland, grassland, scrub, and wetland.
The greater contributions of coarse fragments and texture in alpine vegetation
are consistent with shallow, heterogeneous substrates, whereas the distinct
organic-carbon and carbonate or sodicity contributions in desert niches are
consistent with arid-soil differentiation. These correspondences support the
ecological plausibility of the fitted models, but permutation importance does
not identify causal controls or whether suitability increases or decreases
along a predictor gradient.

Separating climate and soil also determines the interpretation of the future
projection. Temporal change in dual suitability is climate-driven conditional
on a static edaphic gate. Soil restricts where a climatic analogue can be
expressed, but the analysis does not simulate changing soils or
vegetation--soil feedbacks.

\subsection{Binary niche performance versus ranked analogue assignment}

Individual-zone evaluation and multi-ecotype assignment answer different
questions. The former asks whether a focal ecotype can be distinguished from a
sampled background at a fixed threshold. The latter compares 53 independently
fitted suitability values and retains their relative maximum. A cell may
therefore be classified correctly by a focal binary model yet receive another
ecotype as its Top-1 analogue.

The reference-period Top-$k$ results quantify this distinction. Although exact
Top-1 agreement was approximately 60\%, the reference ecotype remained among
the first five analogues across 91\% of evaluated pixels. Most Top-1
disagreements therefore represented comparative ranking among several leading
environmental analogues rather than complete exclusion of the reference
ecotype. Broad-category agreement near 70\% likewise indicates that major
environmental structure was recovered more consistently than boundaries among
closely related ecotypes. The 8.6\% outside the Top-5 nevertheless identifies a
smaller but substantive component not explained by near-rank competition.

Top-$k$ inclusion should not be interpreted as a replacement for Top-1
accuracy. RF vote fractions from independently fitted binary models are not
mutually calibrated posterior class probabilities, and membership in the
Top-3 or Top-5 does not establish equivalent ecological suitability. Ranked
analogues are best treated as a diagnostic of environmental similarity and
assignment sensitivity. They improve interpretation of the hard map but do not
show that a future projection is correct wherever the reference ecotype occurs
within the Top-$k$.

\subsection{Future analogue displacement and threshold-defined novelty}

The primary future comparison used each workflow's modeled reference-period
Top-1 assignment as its baseline. This model-to-model design avoids treating
the reference vegetation raster---with its own reconstruction error---as a
future state. It also separates two ecological-niche outcomes: displacement of
a current analogue within the suitability ranking and absence of any current
analogue above the fixed threshold.

Another current ecotype became Top-1 across much more area than became novel.
Early in the century, most changed cells still retained their former analogue
within the Top-3 or Top-5, indicating that many changes were rank reversals
among comparatively similar niches. By late-century SSP585, however, the
former analogue was outside the Top-5 across almost one-quarter of changed
area. The projected change therefore shifted progressively from modest rank
reordering toward stronger analogue displacement, even though complete
threshold-defined novelty remained less extensive.

Novel niche space reached approximately 0.91--0.96 million~km$^2$ under
late-century SSP585. The simultaneous decline in analogue richness adds a
management-relevant gradient: many areas retained at least one current
analogue but had few alternatives meeting the threshold. For provenance
matching, this distinction separates locations with several candidate source
environments from those with a narrow or absent current analogue set.
Nevertheless, the 0.4 threshold is an operational definition, not an ecological
phase boundary.

Changes in top-ranked assigned area were largest for several managed
agricultural mosaics and for cold, alpine, or subtropical forest ecotypes. The
agricultural-class increases should not be read as forecasts of land-use
expansion or ecosystem replacement; these composite classes carry climatic and
edaphic profiles that can become the highest-ranked environmental analogue.
Likewise, decreases in high-cold meadow, steppe, alpine sparse vegetation, and
subtropical needleleaf niches represent declining relative suitability, not
realized loss. Zone 99 identifies environmental space unmatched by the 53
modeled current ecotypes and does not specify future community composition,
assembly, or persistence.

The matched Top-$k$ maps provide a useful sensitivity bound but not an
independent projection. Their Top-3 and Top-5 assignments restore the observed
reference label when it falls within the corresponding rank set; higher
apparent retention is therefore partly imposed by construction. The fixed
reference-period Top-1 rank analysis is consequently more appropriate for the
primary ecological interpretation.

\subsection{Source-environment variation and assisted migration}

The population-proxy layer combines species occurrences and ecotype niche
models in distinct roles. Occurrences determine which source ecotypes are
represented within a species, whereas the ecotype model supplies the projected
dual-suitability surface. The resulting population proxies are therefore not
independently fitted population distribution models. Population proxies from
the same source ecotype are identical across species, and the species envelope
is the cellwise maximum over the source ecotypes represented for that species.

This construction explains why modest species-level change can conceal strong
source-specific contrasts. The decrease for \textit{L. gmelinii} largely
reflects the cold-temperate and montane-temperate needleleaf niche that
dominates its reference occurrences, whereas the increase for
\textit{Cyclobalanopsis longinux} combines a dominant subtropical evergreen
broadleaf source with several smaller subtropical or tropical sources. For
\textit{Chamaecyparis formosensis}, the approximately 7\% expansion of the
species envelope co-occurred with large contractions in bamboo/scrub,
subtropical needleleaf, and evergreen broadleaf source niches and increases in
two managed mosaic source ecotypes. The latter do not demonstrate adaptation
to cropland or plantations; they reflect the environmental profiles of
composite reference classes.

Retaining source-environment contrasts is relevant to provenance matching
because a species-level envelope treats all included sources as interchangeable.
The cellwise maximum also allows a rare source proxy to extend the envelope as
much as a dominant one, preserving uncommon environmental associations but
making results sensitive to the 10-cell inclusion rule. These proxies can
identify where source-specific suitability differs and where assisted-migration
options merit further testing. They do not establish genetic differentiation,
local adaptation, demographic performance, or establishment success;
application requires common-garden, genomic, demographic, and field evidence.

\subsection{Robustness and limits of inference}

Plain RF and Plain MF RF produced nearly identical predictive metrics, ranked
analogues, projected areas, and feature-importance profiles. Because their
component forests used the same sampled training pool, this agreement indicates
stability to tree number and aggregation rather than robustness to alternative
absence samples. The multiclass RF altered some assignment metrics but retained
the principal reference-period pattern, showing that the main contrast between
binary performance and joint assignment was not unique to one formulation.

Several assumptions constrain inference. Random train--test partitions are not
spatially independent and may overstate transferability to geographically
distinct environments. Independent RF vote fractions can differ in
calibration, yet Top-1 assignment compares them directly. The fixed soil gate
omits soil change and converts continuous edaphic suitability into a hard
constraint; the 0.4 novelty threshold should therefore be examined as a
sensitivity parameter. Correlated predictors redistribute permutation
importance, which cannot supply response direction or causal mechanism.

The reference vegetation map also reflects dispersal history, disturbance,
management, biotic interactions, geomorphic processes, and classification
error that are absent from the predictors. The projections do not model
colonization, demographic persistence, local disappearance, succession, or
ecosystem assembly. They should therefore be interpreted as geographic changes
in modeled environmental suitability for current ecotype analogues. Decisions
about seed transfer, assisted migration, restoration, or genetic conservation
require additional evidence on dispersal, demography, genetic variation,
biotic interactions, disturbance, and establishment.

\section*{Acknowledgments}
% TODO: insert actual funding information (the previous text was template
% boilerplate from an unrelated 2006 grant and must not remain).
[Funding information to be added.]

\section*{Author contributions}
J.J.: conceptualization, methodology, software, formal analysis,
visualization, and writing---original draft. T.W.: supervision, resources,
methodological guidance, and writing---review and editing.

\section*{Competing interests}
The authors declare no competing interests.

\section*{Data and Code Availability}
All analysis code and derived assessment tables are available in the
\href{https://github.com/gingerlolipop/ecoChina2}{ecoChina2 GitHub repository}.
The vegetation map, ClimateAP data, and Harmonized World Soil Database remain
subject to their source-provider access conditions. Paths in the analysis
scripts are local working paths and should be replaced by users' project
directories.

% References
\bibliographystyle{plainnat}
\bibliography{references_corrected_apa}


\section{Supplementary Information}

\setcounter{table}{0}
\renewcommand{\thetable}{S\arabic{table}}
\setcounter{figure}{0}
\renewcommand{\thefigure}{S\arabic{figure}}

\begin{table}[!htbp]
\centering
\caption{Largest changes in Top-1 assigned ecotype-niche area from the
reference-period model map to 2071--2100 SSP585 under Plain RF.}
\label{tab:zone_assignment_change_si}
\begin{threeparttable}
\centering
\scriptsize
\setlength{\tabcolsep}{4pt}
\begin{tabularx}{\textwidth}{@{}rXrrrr@{}}
\toprule
Zone & Short ecotype description &
\shortstack{Reference\\($10^3$ km$^2$)} &
\shortstack{Future\\($10^3$ km$^2$)} &
\shortstack{$\Delta$ area\\($10^3$ km$^2$)} &
Change (\%) \\
\midrule
21 & Grain fields, evergreen orchards, and subtropical economic-tree plantations & 388.7 & 1332.7 & +944.0 & +242.8 \\
27 & Grain fields and deciduous orchards & 488.6 & 835.5 & +346.9 & +71.0 \\
18 & Subalpine broadleaf deciduous scrub & 116.9 & 437.4 & +320.5 & +274.1 \\
53 & Grain fields, tropical evergreen orchards, and economic-tree plantations & 76.7 & 255.4 & +178.8 & +233.2 \\
46 & Subalpine sclerophyllous broadleaf evergreen scrub & 129.6 & 307.1 & +177.6 & +137.1 \\
12 & Grain fields and cold-tolerant economic crops & 213.4 & 356.8 & +143.4 & +67.2 \\
\addlinespace
16 & \textit{Kobresia} and forb high-cold meadow & 503.6 & 169.7 & -333.9 & -66.3 \\
30 & Subtropical needleleaf forest & 351.2 & 52.7 & -298.5 & -85.0 \\
31 & Grass and \textit{Carex} high-cold steppe & 524.8 & 276.8 & -248.0 & -47.3 \\
47 & Grain fields, orchards, and economic-tree plantations & 266.3 & 33.8 & -232.5 & -87.3 \\
20 & Alpine sparse vegetation & 293.8 & 63.2 & -230.6 & -78.5 \\
13 & Temperate tufted-grass steppe & 338.5 & 109.3 & -229.2 & -67.7 \\
\bottomrule
\end{tabularx}
\begin{tablenotes}[flushleft]
\footnotesize
\item[] \textbf{Notes.} Areas are weighted by cell-specific surface area.
Values describe the area in which each ecotype niche was ranked first after
comparison among all 53 models; they are not estimates of realized vegetation
extent and exclude cells where the ecotype was suitable but ranked lower.
\end{tablenotes}
\end{threeparttable}
\end{table}

\begin{figure}[!htbp]
  \centering
  \projectfigure[width=0.98\linewidth]{Figure_var_3_zone_level_climate_soil_metrics.png}
  \caption{Held-out climate and topsoil model performance by ecotype.}
  \label{fig:zone_performance_si}
\end{figure}

\begin{figure}[!htbp]
  \centering
  \projectfigure[width=0.98\linewidth]{Figure_var_4_reference_topk_confusion_flows_common_scale.png}
  \caption{Largest residual reference-period ecotype confusion flows under
  Top-1, Top-3, and Top-5 agreement. All panels use a common vertical scale;
  for $k>1$, a pixel agrees when its reference ecotype occurs within the first
  $k$ ranks.}
  \label{fig:topk_confusion_si}
\end{figure}

\begin{figure}[!htbp]
  \centering
  \projectfigure[width=0.94\linewidth]{Figure_S_analogue_scarcity_by_minimum_count.png}
  \caption{Projected richness of suitable current-ecotype analogues. The
  cumulative criteria are zero, fewer than three, and fewer than five ecotypes
  with dual suitability of at least 0.4. Only the zero-analogue category is
  novel niche space.}
  \label{fig:analogue_scarcity_si}
\end{figure}

\begin{figure}[!htbp]
  \centering
  \projectfigure[width=0.98\linewidth]{Figure_var_7_matched_topk_ecotype_niche_change_rf_var.png}
  \caption{Matched Top-$k$ sensitivity analysis for Plain RF. Reference and
  future maps use the same rank criterion. For Top-3 and Top-5, the observed
  reference ecotype is restored when it occurs among the first $k$ ranks;
  these panels are reference-conditioned diagnostics rather than independent
  future maps.}
  \label{fig:matched_topk_rf_si}
\end{figure}

\begin{figure}[!htbp]
  \centering
  \projectfigure[width=0.98\linewidth]{Figure_var_7_matched_topk_ecotype_niche_change_mf_var.png}
  \caption{Matched Top-$k$ sensitivity analysis for Plain MF RF, defined as in
  Supplementary Fig.~\ref{fig:matched_topk_rf_si}.}
  \label{fig:matched_topk_mf_si}
\end{figure}

\begin{figure}[!htbp]
  \centering
  \begin{subfigure}{0.98\linewidth}
    \centering
    \projectfigure[width=\linewidth]{Figure_FI_reader_climate_period_composition.png}
    \caption{Annual and seasonal composition of climatic importance.}
  \end{subfigure}

  \vspace{0.5em}

  \begin{subfigure}{0.98\linewidth}
    \centering
    \projectfigure[width=\linewidth]{Figure_FI_reader_rf_vs_mf_importance_agreement.png}
    \caption{Agreement of feature-importance profiles between RF workflows.}
  \end{subfigure}
  \caption{Supplementary diagnostics for feature-importance interpretation.}
  \label{fig:importance_diagnostics_si}
\end{figure}

\begin{figure}[!htbp]
  \centering
  \projectfigure[width=0.96\linewidth]{Figure_var_10b_reference_map_overall_metrics.png}
  \caption{Reference-period agreement under the two individual-zone workflows
  and the multiclass RF. Error bars on macro ecotype-level metrics are SE
  across ecotypes; map-wide percentages are deterministic summaries of the
  common raster mask.}
  \label{fig:robustness_multiclass_si}
\end{figure}

\begin{figure}[!htbp]
  \centering
  \projectfigure[width=0.98\linewidth]{Figure_var_5_future_maps_mf_var.png}
  \caption{Projected Top-1 ecotype-niche assignments from Plain MF RF under
  SSP245 and SSP585. Zone 99 denotes novel niche space.}
  \label{fig:future_maps_mf_si}
\end{figure}



\end{document}
